\section{安全与成本分析}

\subsection{防篡改证据与可追踪性}

TrustAuditFlow 使用两类哈希结构实现可追踪性。首先，分片哈希形成 Merkle 根。一旦根被锚定到链上，后续对任一分片的修改都会改变对应叶子，并使 Merkle 路径失效。其次，审计事件按哈希链排序。如果攻击者在锚点提交后删除、插入或修改链下事件，重新计算得到的哈希链状态将无法匹配已锚定状态。

这并不意味着链下证据自动可用。证据存储仍必须保留原始审计材料。这里的贡献在于，链能够检测已保存证据与已提交时间线之间的不一致。

\subsection{委员会安全边界}

委员会路径降低常规确认成本，但也引入选择风险。如果攻击者控制目标对象的委员会多数，委员会可能延迟或批准虚假审计摘要。TrustAuditFlow 通过将委员会仅作为轻量路径来处理该风险：委员会候选必须通过信任阈值，委员会成员应轮换，高风险事件使用更大委员会或全节点确认，任何法定人数失败都会触发回退。

因此，委员会路径在有界安全范围内使用。它是在显式回退规则下针对常规事件的可配置优化，而完整 BLoP 可信节点集合仍是争议事件或高风险事件的高保障路径。

\subsection{链上成本}

链上成本来自载荷数量、载荷大小和共识确认。批量锚定将常规事件的记录数量从逐事件记录降低为逐批次锚点。委员会确认在风险感知路径下降低常规事件的投票节点数量。纠删码相比朴素复制降低了保持证据可恢复所需的冗余。通过这些机制，系统将常规审计流量与高保障全节点确认分离。

\subsection{与 PDP 和 PoR 的关系}

TrustAuditFlow 不替代 PDP 或 PoR。PDP/PoR 协议证明数据持有或可恢复性；TrustAuditFlow 组织挑战策略、证明摘要、异常状态、恢复结果和问责记录如何被锚定与确认。基于哈希的校验、公开 PDP/PoR 校验和零知识证明模块均可通过相同的审计证据接口集成。

